%%%%%%%%%%%%
%
% $Autor: Asena Yüce 7024749, Carl-Lewis Maschke 7024640, Nina Braams 7024847, Sven Neumann 7025068, Valeria Kornichenkova 7024769 $
% $Datum: 2019-03-05 08:03:15Z $
% $Pfad: Maintenance.tex $
% $Version: 4250 $
% !TeX spellcheck = en_GB/de_DE
% !TeX encoding = utf8
% !TeX root = manual 
% !TeX TXS-program:bibliography = txs:///biber
%
%%%%%%%%%%%%

\chapter{Wartung}
	\textbf{Allgemeiner Hinweis:}\\
	Für jegliche Wartungsarbeiten muss die Stromzufuhr unterbrochen werden. Die Stromzufuhr darf erst nach Beendigung der Wartungsarbeiten eingeschaltet werden 
	\textbf{Wartungshinweise:}
	\begin{enumerate}
		\item Nach längeren Stillstandszeiten sollte das Wasser im Blumentopf ausgetauscht werden. Hierzu wird der Deckel mitsamt Acrylrohr abgenommen. Das abgestandene Wasser wird entleert, und der Blumentopf kann bei Bedarf gründlich gereinigt werden. Anschließend wird frisches Wasser bis zur maximalen Füllhöhe eingefüllt. Zum Abschluss wird der Deckel mit dem Acrylrohr wieder aufgesetzt.
		\item Zur Reinigung des Messrohrs wird zunächst der Deckel abgenommen und das Messrohr vom Deckel abgeschraubt. Danach kann der Verschluss entfernt und der Mess-Schwimmer entnommen werden. Messrohr und Verschluss werden gründlich gereinigt. Nach der Reinigung wird der Verschluss wieder auf das Messrohr aufgesetzt. Der Mess-Schwimmer wird von der offenen Seite her in das Messrohr eingeführt, das anschließend wieder am Deckel montiert wird.
		\item Für die Reinigung des Ultraschallsensors wird ebenfalls der Deckel abgenommen und das Messrohr abgeschraubt. Der Sensor selbst darf dabei nicht vom Deckel entfernt werden. Die Reinigung erfolgt ausschließlich mit einem Mikrofasertuch in sanften Bewegungen. Ein direkter Hautkontakt mit dem Sensor ist unbedingt zu vermeiden. Nach der Reinigung wird das Messrohr wieder am Deckel befestigt und der Deckel abschließend auf dem Blumentopf platziert.
		\item Soll der Mess-Schwimmer gereinigt oder ersetzt werden, wird zunächst der Deckel abgenommen. Anschließend wird der Verschluss des Messrohrs abgeschraubt und der Schwimmer entnommen. Nach der Reinigung oder dem Austausch wird der Mess-Schwimmer wieder in das Messrohr eingesetzt und der Verschluss ordnungsgemäß aufgeschraubt.
	\end{enumerate}
	 


